\bookchapter{第三章：使手機接“生命之線”——網絡}

\chapterintro{無網，智能手機則“智障”其半。網絡乃現代手機之生命之線，解之、用之、護之，乃每個手機用戶之必修課。}

\booksection{〇、聯網不迷——分清Wi-Fi與流量之“路”與“車”}

衆人之手機屏幕上方同時亮着Wi-Fi與信號之圖標，然於二者之本區別卻一頭霧水。最經典之誤即：\textbf{連着Wi-Fi卻刷不出頁}——蓋因路由器斷網，而手機未自動切迴流量。欲徹底明此問題，須建一“路與車”之認知之模。

\begin{tcbinfo}{核心之喻：路與車}
視互聯網如一城。\textbf{“路”}乃汝手機至互聯網之間之通道，\textbf{“車”}乃行於路上之數據。Wi-Fi與流量乃兩條不同之路，然數據（車）之本身則一也。

\end{tcbinfo}

\booksubsection{本質之別：本地通道 vs 遠程高速}

Wi-Fi本質上乃一種\textbf{局域網之技}，猶汝家樓下無線鋪設之路，司將數據自路由器傳至手機。然若寬帶欠費或光纜斷，路雖猶在，外之“高速公路”卻已斷，自無處可往。而網絡流量（移動數據）則經運營商之4G/5G基站傳輸，相當於汝直開於運營商之高速專線，有信號即直連互聯網。

\booksubsection{費用與場：固定月租 vs 按量計費}

Wi-Fi通常爲一次性之投（寬帶月費），但電費不斷，全家無限用，甚宜下電影、玩遊戲等大流量之場。而流量則按用量計費，超套餐或降速或額外扣費，更宜於地鐵、出差等移動之場處理輕量之事。

\begin{tcbpractice}{來一試：檢汝之流量套餐}
開手機之 [設置] > [移動網絡] > [流量管理]，觀今月尚餘幾何流量、何App最耗流量。心中有數，出門方不慌。

\end{tcbpractice}

\booksubsection{安之慮：公共Wi-Fi乃隱之“釣魚場”}

於安之上，移動數據通常勝於公共Wi-Fi。移動數據由運營商加密，而公共Wi-Fi實易爲黑客搭爲“僞熱點”截獲未加密之流量。然此處須分清：於HTTPS全普及之今，攻擊者能截獲者乃\textbf{汝所訪之網站（域名）}，以及\textbf{未加密之HTTP內容}；而微信消息、銀行App之登錄密碼等經HTTPS加密之數據，攻擊者\textbf{不能直讀}——除非施更復雜之SSL剝離攻擊且用戶略之證書之警。然即便如此，於外行支付之事，直用手機流量仍爲最穩之擇。

\textbf{安全提醒}\par

於機場、商場、咖啡廳等公共之所，勿連名號可疑之免費Wi-Fi（如“Free\_WiFi\_2024”）。縱連之，亦絕勿於其上登銀行App或輸密碼。

\booksubsection{手機端之判：解狀態欄之“暗語”}

何以知手機此刻用何聯網？以華爲手機爲例，狀態欄中運營商信號旁或WLAN信號旁現\textbf{雙向箭之圖標}，即表正用對應之方上網；若二者同時亮，則網絡加速生效，手機正同時用兩路。

iPhone用戶則更直觀：Wi-Fi圖標亮即在用Wi-Fi，不亮時觀右上角之信號格即在用流量。部分安卓手機亦於狀態欄顯“5G”或“4G”之字，徑告汝今之移動網絡之類。

\booksection{一、移動網絡（4G/5G）與流量}

移動網絡乃汝隨時隨地聯網之基。自2G至5G，速度增上千倍，然其理實同。

\booksubsection{網絡之制：}

\begin{itemize}
  \item \textbf{4G LTE}：今覆蓋最廣，速度足用（理論上最高百Mbps+），耗電相對低。
  \item \textbf{5G}：速快（理論上最高10Gbps），然覆蓋猶在完善，耗電高。宜下大文件、觀高清視頻。
  \item \textbf{5G SA vs NSA}：SA爲獨立組網，NSA乃依託4G基站。SA之體驗更佳，然須運營商與手機皆支持。
\end{itemize}

\booksubsection{流量之管：}

\begin{itemize}
  \item \textbf{流量何物}：即手機上網所生之數據量。觀視頻、刷圖最費流量，文字與語音相對省。
  \item \textbf{省流之小貼}：於設置中開"省流量模式"，Wi-Fi下再更App或觀高清視頻。
  \item \textbf{信號格$\neq$網速}：信號強僅表連接佳，網速猶視基站之忙閒、汝距基站之遠近。
\end{itemize}

\textbf{重要之警：出國前請閉【數據漫遊】！}\par

於國外用移動數據或生天價之賬單！出國前務於設置中閉數據漫遊，或先購當地流量卡/國際漫遊套餐。

\booksection{二、Wi-Fi：連接、分享與安}

Wi-Fi乃家與辦公室聯網之主將，速快、不費流量。然Wi-Fi實亦不斷進化，不同之標差甚大。

\booksubsection{連Wi-Fi：}

往 [設置] > [WLAN] 擇網絡輸密碼。今諸多手機支持掃碼連接，更爲便。

\booksubsection{分享密碼：}

安卓與iOS皆能使已連之手機速分享密碼於友。安卓生二維碼，iOS可靠近自動分享。

\booksubsection{Wi-Fi之標演進：}

Wi-Fi之標由IEEE制定，每代皆顯著提升：

\begin{center}
\begin{tabularx}{\linewidth}{|X|X|X|X|X|}
\hline
\textbf{標準} & \textbf{協議} & \textbf{頻段} & \textbf{最高速率} & \textbf{特點} \tabularnewline
\hline
Wi-Fi 4 & 802.11n & 2.4/5GHz & 600Mbps & 老舊之標，覆蓋尚可 \tabularnewline
Wi-Fi 5 & 802.11ac & 5GHz & 3.5Gbps & 今最普及，速度足用 \tabularnewline
Wi-Fi 6 & 802.11ax & 2.4/5GHz & 9.6Gbps & 多設備之能強，延遲低 \tabularnewline
Wi-Fi 6E & 802.11ax & 2.4/5/6GHz & 9.6Gbps & 增6GHz頻段，不擁堵 \tabularnewline
Wi-Fi 7 & 802.11be & 2.4/5/6GHz & 46Gbps & 極速傳輸，多鏈路聚合 \tabularnewline
\hline
\end{tabularx}
\end{center}

\begin{itemize}
  \item \textbf{Wi-Fi 6之關鍵技}：OFDMA（正交頻分多址）使多設備同傳，MU-MIMO（多用戶多入多出）提升多設備併發之效，TWT（目標喚時）使設備更省電。
  \item \textbf{Wi-Fi 6E}：於Wi-Fi 6之基增6GHz頻段，供更多信道，減干擾。須手機與路由器皆支持。
  \item \textbf{Wi-Fi 7}：引入MLO（多鏈路操作），可同時使用多個頻段傳輸數據，大幅降低延遲。320MHz超寬信道使速率飛躍。驍龍8 Elite Gen5、天璣9400等旗艦芯片平臺均已支持Wi-Fi 7，小米17、華爲Mate 80、三星S26等機型可使用Wi-Fi 7網絡（需配合Wi-Fi 7路由器）。
  \item \textbf{2.4GHz vs 5GHz}：2.4GHz穿牆佳而速慢、干擾多；5GHz速快而穿牆弱。Wi-Fi 6E/7之6GHz頻段速最快，而穿牆最弱，宜同一室中用。
\end{itemize}

\booksubsection{公共Wi-Fi之險：}

公共Wi-Fi或爲黑客竊聽數據或設釣魚熱點。守則：

\begin{itemize}
  \item 避登銀行賬戶、輸密碼。
  \item 最佳用VPN加密連接。
  \item 勿自動連陌生Wi-Fi。
  \item 用HTTPS網站（地址欄有鎖之圖標）。
\end{itemize}

\booksection{三、藍牙：短距無線通訊之萬能手}

藍牙或爲汝手機上用頻最高之無線技之一——無線耳機、智能手錶、車載系統、鍵鼠，皆賴其連。

\booksubsection{技之理：}

藍牙工作於2.4GHz ISM頻段，用跳頻擴頻之技（FHSS）將數據散於七十九頻道上速切傳，有效避干擾。藍牙之通訊距分三級：Class 1約百米、Class 2約十米（手機常用）、Class 3約一米。

\booksubsection{版本之演進：}

\begin{itemize}
  \item \textbf{藍牙4.0}：引入低功耗藍牙（BLE），使手環、心率帶此類小電池之設備亦能無線連，續航自天變月。
  \item \textbf{藍牙4.2}：提升傳輸之速與隱私之護，支持IPv6直連。
  \item \textbf{藍牙5.0}：傳輸之速倍增（2Mbps），傳輸之距四倍（理論三百米），廣播之能八倍。支持雙音頻（同時連兩副耳機）。
  \item \textbf{藍牙5.1}：引入方向尋址之能（AoA/AoD），定位之精可達釐米級。
  \item \textbf{藍牙5.2}：LE Audio新標，支持LC3之編碼（音質更佳、功耗更低）、多流音頻、廣播音頻。
  \item \textbf{藍牙5.3/5.4}：進一步優化功耗與連之穩，增強週期性廣播。
\end{itemize}

\textbf{LE Audio之普及況}：2025-2026年，LE Audio已進入快速普及階段。主流品牌之新款真無線耳機（如華爲FreeBuds Pro 5、小米Buds 6 Pro、三星Galaxy Buds3 Pro等）均已支持LE Audio與LC3編碼，提供更低延遲與更好音質。購買新耳機時，建議優先選擇支持LE Audio之產品。

\booksubsection{應用之場：}

\begin{itemize}
  \item \textbf{音頻之傳}：無線耳機、車載之響。藍牙5.2之LE Audio使音質與延遲皆有質之飛躍。
  \item \textbf{數據之傳}：近距傳文件、照。雖速不如Wi-Fi Direct，然勝於簡。
  \item \textbf{設備之配}：智能手錶、手環、鍵鼠、遊戲手柄。
  \item \textbf{位置之服}：藍牙信標（Beacon）用於室內導航、物之追蹤（如蘋果AirTag）。
  \item \textbf{智能之家}：藍牙Mesh網控燈泡、門鎖等。
\end{itemize}

\booksection{四、紅外遙控：手機化萬能遙控}

紅外遙控乃最古之無線控技之一，然於手機上猶爲實用。部分安卓手機（如小米、華爲部分機型）內置紅外發射器，可遙控電視、空調、機頂盒等家電。

\booksubsection{工作之理：}

紅外遙控經紅外發光二極管發特定頻（通常38kHz）之紅外光脈衝信號，家電之紅外接收器解此脈衝而執對應之操。不同之品牌用不同之編碼之協議，常見者有NEC、RC5、RC6、SONY等。手機中之紅外App內置各品牌之編碼庫，擇品牌與型號即能配。

\booksubsection{局之限：}

\begin{itemize}
  \item \textbf{方向之性}：紅外乃直線傳，須對設備發之，不能隔牆。
  \item \textbf{距之短}：有效之距一般五至八米。
  \item \textbf{蘋果不支}：iPhone無紅外發射器，不能用此功。
\end{itemize}

\booksection{五、NFC：碰一碰之法}

NFC（Near Field Communication，近場通訊）乃一種極短距（十釐米以內）之高頻無線通訊之技，頻十三點五六MHz，傳輸之速最高424Kbps。

\booksubsection{通訊之理與模：}

NFC本於RFID技而變，支持三作之模：

\begin{itemize}
  \item \textbf{讀寫之模}：手機讀NFC標籤中之信，如讀NFC貼紙而觸動作。
  \item \textbf{點對點之模}：兩臺NFC設備互傳數據，如Android Beam（已停用）。
  \item \textbf{卡模擬之模}：手機模擬一智能卡，此乃最常用之模——公交卡、門禁卡、銀行卡皆此模也。
\end{itemize}

\booksubsection{安之機：}

\begin{itemize}
  \item \textbf{超短之距}：十釐米內之通訊距本身即物理之安，遠竊幾乎不能。
  \item \textbf{安之芯}：手機內置獨立之安之芯（SE）或用HCE（主機卡模擬），支付之鑰存於安之區，即手機被root亦不能提。
  \item \textbf{指紋/面容之驗}：支付時鬚生物之識以認，即手機失亦不能盜刷。
  \item \textbf{令牌之化}：Apple Pay、華爲Pay等不傳真卡號，乃用一次之令牌，即被截亦不能複用。
\end{itemize}

\booksubsection{典型之用：}

\begin{itemize}
  \item \textbf{公交卡/門禁卡}：將實體卡之信寫於手機，刷手機即行。
  \item \textbf{移動之付}：Apple Pay、華爲Pay等，將銀行卡綁於手機，碰一碰即付。
  \item \textbf{設備之配}：NFC碰一碰速配藍牙耳機、音箱。
  \item \textbf{身份之驗}：NFC身份證之讀、電子駕照之驗。
\end{itemize}

\booksection{六、個人熱點：將手機變Wi-Fi路由器}

己流量多時，可開熱點以予平板、電腦用。

\begin{itemize}
  \item \textbf{開啓之法}：設置 > 個人熱點/網絡共享，開啓之。
  \item \textbf{注意之事}：甚耗電與流量，宜設一複雜之密碼，並限連設備之數。
  \item \textbf{USB之共}：以數據線連電腦，較無線之點更穩、更快。
\end{itemize}

\begin{tcbpractice}{動手一試：設並管熱點}
\begin{enumerate}
  \item 開 [設置]，尋"個人熱點"或"網絡共享"，開啓之，設好名與密碼。
  \item 用另一設備（如電腦）搜而連此Wi-Fi。
  \item 於手機端觀之，能否理已連之設備，或踢去不識之設備。
\end{enumerate}

\end{tcbpractice}
